\section{Noise Management Techniques in RNS}
\label{sec:gadgets}
A problem that arises in several areas of \gls{fhe} is that of multiplying a ciphertext by some large, known factor $X\in R_Q$ while keeping the error low. Performed naively, the noise scales together with the message:
\begin{equation}
    (a\cdot s + \Delta m + \epsilon) \cdot X = aX\cdot s + \Delta Xm + X\epsilon.
\end{equation}

\textcolor{purple}{
Several core operations require computing the product of a ciphertext element with noisy auxiliary material; performed naively, this scales the noise by a factor as large as $q$. This section covers the standard techniques for controlling this growth: decomposition (BV), modulus extension (GHS), and their combination (Hybrid).
}

These techniques typically occur in the literature in the context of key switching, but we remove them entirely from that context and consider only the techniques themselves.

\subsection{Brakerski-Vaikuntanathan}
Brakerski and Vaikuntanathan~\cite{cryptoeprint:2011/344} first proposed decomposing ciphertext elements into small digits as a means of controlling the noise growth of this operation. The idea has since been generalized into the idea of \emph{gadget decomposition} and applied in a variety of forms, including digit and CRT (RNS) decomposition. Following the convention of~\cite{cryptoeprint:2021/204, cryptoeprint:2022/915}, we refer to this family of decomposition-based techniques collectively as \emph{BV}.

% \textcolor{purple}{\textit{Decomposition}; a technique used to represent a number in a smaller base (radix) while preserving its value.} % Taken from https://fhetextbook.github.io/Decomposition.html#decomposition

\begin{definition}[Digit Decomposition]
    \label{ex:radix_gadget}
    takes an integer $\gamma\in\mathbb{Z}_q$ (for a modulus $q\geq 2$), base $\beta$ and \textit{decomposition parameter} $\ell$, and maps $\gamma$ to its $\ell$ most significant \textit{signed} base-$\beta$ digits $\gamma_i \in [-\beta/2, \beta/2)$:
    \begin{equation}
        \gamma = \sum_{i=0}^{\ell-1}\gamma_i\frac{q}{\beta^{i+1}} + \epsilon, \qquad |\epsilon| \leq \frac{q}{2\beta^{\ell}},
    \end{equation}
    where the remaining, least significant digits are simply dropped; keeping all digits ($\beta^\ell \geq q$) makes the decomposition exact ($\epsilon = 0$). The \textit{decomposition of} $\mathit{\gamma}$ is the ordered list of the $\ell$ digits and is denoted $\mathsf{Decomp}(\gamma) = (\gamma_0, \gamma_1, \dots, \gamma_{\ell-1})$.
\end{definition}

\begin{remark}
    When $q$ is not a power of $\beta$ (e.g.\ the prime moduli of \autoref{sec:rns}) the entries $q/\beta^{i+1}$ are rounded to the nearest integer; the rounding error is absorbed into $\epsilon$.
\end{remark}

%%
%% Highlight what's imporant for generalizing
\textcolor{purple}{
By construction, digit decomposition satisfies $\langle \mathsf{Decomp}(\gamma),\, g \rangle = \gamma$ for the fixed vector $g = (q/\beta, \dots, q/\beta^{\ell})$. The particular form of $g$ is unimportant; what matters is that the identity holds while the entries of $\mathsf{Decomp}(\gamma)$ remain small, since then, scaling $g$ by any $x \in \mathbb{Z}_q$, \[ \langle \mathsf{Decomp}(\gamma),\, x \cdot g \rangle = x \cdot \gamma \pmod q, \] the product $x \cdot \gamma$ is computed with $\gamma$ entering only through its small digits. The gadget property abstracts exactly this: it requires nothing of the maps $D(\gamma) = \mathsf{Decomp}(\gamma)$ and $P(x) = x \cdot g$ beyond that their inner product (approximately) recovers $x \cdot \gamma$ and that the entries of $D(\gamma)$ are small.
}

\begin{definition}[Gadget Property.]
    \label{def:gadget}
    Two mappings $P(x), D(y) \in R_q^\ell$ are said to have the gadget property with \textit{quality} (digit bound) $\beta$ and \textit{precision} $\epsilon$ iff. they satisfy \eqref{def:gadget-property} for any pair of integers $(x, y) \in \mathbb{Z}_q^2$.
    \begin{equation}
        \label{def:gadget-property}
        \langle P(x), D(y)\rangle = x\cdot (y + \epsilon_y) \pmod{q},
        \qquad |\epsilon_y| \leq \epsilon, \quad \|D(y)\|_\infty \leq \beta
    \end{equation}
    If $\epsilon = 0$ then $P(x), D(y)$ are said to have the \textit{exact} gadget property.
\end{definition}

\begin{definition}[Gadget Decomposition]
    generalizes the notion of digit decomposition by replacing the base $\beta$ with a \textit{gadget vector} $g = (g_0, g_1, \dots, g_{\ell-1})$ containing $\ell$ elements, such that $D(\gamma) = \mathsf{Decomp}^g(\gamma)$ and $P(x) = x\cdot g$ satisfy Def.~\ref{def:gadget}.

    \begin{equation*}
        \label{def:gadget-decomp}
        \gamma + \epsilon_\gamma = \sum_{i=0}^{\ell - 1} \gamma_i g_i \quad\iff\quad \langle \mathsf{Decomp}^g(\gamma), g\rangle = \gamma + \epsilon_\gamma
    \end{equation*}
\end{definition}

All of the above extends coefficient-wise to polynomials $\gamma \in R_q$: the digits become polynomials with small coefficients, and the bounds $\beta$ and $\epsilon$ apply to $\|\cdot\|_\infty$.

\subsubsection{RNS Instantiation}
\label{sec:gadgets:rns}
Digit decomposition is not directly compatible with RNS: the digits of $\gamma$ are defined by its canonical value, which is not available limb-wise. Extracting them requires reconstructing the coefficients --- an inverse NTT followed by an inverse CRT --- reintroducing the multi-precision detour that RNS is meant to avoid (\autoref{sec:rns-base-conversion}).

The CRT structure itself, however, is a gadget. The first RNS instantiation of BV~\cite{cryptoeprint:2016/510} decomposes into residues:
\begin{equation}
    \label{eq:crt-gadget}
    \gadgetformat{D}{CRT}(\gamma) = \left([\gamma]_{q_0}, \dots, [\gamma]_{q_{k-1}}\right), \qquad
    \gadgetformat{P}{CRT}(x) = \left(x\cdot e_0, \dots, x\cdot e_{k-1}\right), \qquad
    e_j = \left[\left(\frac{Q}{q_j}\right)^{-1}\right]_{q_j}\cdot\frac{Q}{q_j},
\end{equation}
where the $e_j$ are the CRT idempotents: $e_j \equiv 1 \pmod{q_j}$ and $e_j \equiv 0 \pmod{q_{j'}}$ for $j' \neq j$, so that $\sum_j [\gamma]_{q_j}\, e_j \equiv \gamma \pmod{Q}$ --- the exact gadget property with $\ell = k$. Both maps are computed natively in RNS: $\gadgetformat{D}{CRT}$ simply reads off the limbs, and $\gadgetformat{P}{CRT}$ scales them. The drawback is the digit bound $\beta = \max_j q_j/2$, a full machine word, which enters the noise bounds of \autoref{sec:tfhe} linearly. \cite{cryptoeprint:2018/117} (HPS) further optimized this construction for key switching using precomputed factors, as remarked in \autoref{sec:rns-base-conversion}.

To shrink the digits, the CRT gadget is composed with digit decomposition applied \textit{per residue}: each limb $[\gamma]_{q_j}$ is an ordinary word-sized integer, so it can be decomposed into $\ell$ signed base-$\beta$ digits limb-wise, without reconstruction. The composed gadget
\begin{equation}
    \label{eq:bv-gadget}
    \bv{D}(\gamma) = \left(\mathsf{Decomp}\left([\gamma]_{q_j}\right)\right)_{j=0}^{k-1}, \qquad
    \bv{P}(x) = \left(x\cdot \beta^i e_j\right)_{j,i}
\end{equation}
has $k\ell$ digits bounded by $\beta/2$, and is exact whenever $\beta^\ell \geq \max_j q_j$. This is the closest direct translation of TFHE-style decomposition to RNS, and one of the two variants implemented in this thesis (\autoref{chp:methodology}).

\subsection{Gentry-Halevi-Shoup}
\label{sec:gadgets:ghs}
The GHS method~\cite{cryptoeprint:2012/099} controls noise without any decomposition ($\ell = 1$). Instead, the modulus is temporarily extended from $Q$ to $QP$ for an auxiliary modulus $P$ (\autoref{sec:rns-base-conversion}), and the auxiliary material is pre-scaled by $P$:
\begin{equation}
    \label{eq:ghs-gadget}
    \gadgetformat{D}{GHS}(\gamma) = \mathsf{ModUp}(\gamma) \in R_{QP}, \qquad
    \gadgetformat{P}{GHS}(x) = [P\cdot x]_{QP},
\end{equation}
where $\mathsf{ModUp}$ is the fast base extension of Algorithm~\ref{alg:fast_base_extension}. The product $\gadgetformat{D}{GHS}(\gamma)\cdot\gadgetformat{P}{GHS}(x) = P\cdot x\gamma \pmod{QP}$ is brought back to $Q$ with $\mathsf{ModDown}$ \eqref{eq:approxmoddown}, dividing by $P$. The single ``digit'' $\gamma$ remains as large as $Q$, but any noise attached to the pre-scaled material is divided by $P \approx Q$ on the way down, shrinking to insignificance; the price is a small additive error $\epsilon > 0$ (the FastBConv overflow and the ModDown rounding) and auxiliary material living in the larger ring $R_{QP}$. GHS is well established for key switching; its use in the external product (\autoref{sec:tfhe}) is not, and is developed in this thesis (\autoref{chp:methodology}).

\subsection{Hybrid}
BV and GHS are endpoints of a spectrum. Hybrid key switching~\cite{cryptoeprint:2021/204} groups the $k$ limbs into $\mathsf{dnum}$ digits of $\alpha = \lceil k/\mathsf{dnum}\rceil$ limbs each --- a coarser CRT gadget --- and applies the GHS modulus extension to control the remaining per-digit noise, requiring only $P \gtrsim Q^{1/\mathsf{dnum}}$. Setting $\mathsf{dnum} = k$ (no extension needed) recovers BV, while $\mathsf{dnum} = 1$ recovers GHS; the ``hybrid'' implementation of this thesis uses $\mathsf{dnum} = 1$ and is therefore effectively GHS.
